\section{Guenon on the Young Evola}

Two letters from Rene Guenon to Guido de Giorgio regarding Julius Evola. Although Evola translated and promoted Guenon's works in Italy, Guenon had some reservations about Evola's approach. Unfortunately, because of that association, certain political views have been attributed to Guenon that he never personally held.

\begin{quotex}
He is a fool because he does not know what everyone knows; he is a sage because he knows what everyone does not know.
\flright{\textsc{Fulvio Mocco}, \emph{I Tarocchi esoterici}}
\end{quotex}


\textbf{Rene Guenon} was described by \textbf{Julius Evola} on many occasions as the `Master of the 20\textsuperscript{th} century'. By the time Evola began his public writing career, Guenon had already had more than 20 years of experience in Tradition, including initiations into Hermetism, Sufism, and Vedanta. Guenon had no interest in Western philosophy as such. In his trilogy Man and his becoming, Symbolism of the Cross, and Multiple States of Being, he expounded Traditional metaphysics, based primarily on the Vedanta and Sufism. Therein is revealed an understanding of non-duality, the ultimate reality, in terms suitable for a Westerner.

When we discuss Evola, we strive to do so, in the light of Guenon. When there is a discrepancy in formulations, Guenon's position needs to be evaluated very carefully before discarding it. We understand the appeal of Evola; he provides a virile spirituality in a world of effeminate churches and clergy. Yet we must urge those who come to Tradition only through Evola to also delve deeply into Guenon.

We provide here two letters from Guenon to \textbf{Guido de Giorgio} regarding Evola, when Evola was under 30 years old. De Giorgio, an initiate into Sufism and a colleague of Guenon, introduced Evola to Tradition in the Ur group. As Guenon notes, Evola's misunderstandings flow from his philosophical position, which is ultimately incompatible with metaphysics. Evola's protests don't ring true, since, as we saw in a recent post regarding Mircea Eliade\footnote{\url{https://gornahoor.net/?p=4303}}, the former's philosophy was integral to his work.

For Evola, the ultimate was the Absolute Individual, which is not the same as non-dual metaphysics. He also took Bachofen's notion of Solar and Lunar cultures too far, applying them to various traditions. Now a legitimate tradition is based on principle, so the solar/lunar distinction is beside the point. He was also attached to Weininger's theory of the Absolute Male and Absolute Female. That may be interesting on the psychological level, but is misleading on a spiritual or transcendental level. Male and Female form a polarity, united by a third term (see Guenon's Great Triad). They are also relative. For example, a knight is masculine in relation to the King's subjects, but feminine in relation to the king. To absolutize one pole would be like making Yang the supreme principle rather than the Tao.

In the letters from Guenon to de Giorgio, we see Guenon's early frustration with the youthful Evola. They speak for themselves to those who understand what is at stake. The first letter is from 20 November 1925.

\begin{quotex}
Thank you for the issue of Bilychnis that you sent me; since then, I received from Evola himself a complete packet of other journals containing some of his articles. In acknowledging its reception, I told him that I had some serious reservations to make on his point of view that appeared to me too philosophical. He wrote me a somewhat long, rather confused, letter, in which he protested that the philosophical form he uses is for him only a simple means of exposition that does not affect his doctrine itself. I believe none of it. I persist in thinking that he is really too imbued with philosophy, and especially with German philosophy. In an article published in the journal Ultra, he alluded to me in a note, concerning East and West, in terms that prove that he did not understand very much about what I expounded; all the same, he goes as far to describe me as a `rationalist', which is really ridiculous (considering that it concerns a book where I expressly asserted the falsity of rationalism!), and which clearly shows that he is among those who cannot get rid of philosophical labels and who feel the need to apply them wrongly everywhere. He informed me of his intention to write an article on Man and his becoming; I ask myself what that could be like; we'll see eventually... Since we are speaking of Evola, it is yet necessary that I tell you that he was hurt by Reghini's criticisms of him, even though in a very moderated form. He must be rather vain and would like to have nothing but eulogies; it is true that he is very young. Vulliaud, who does not have the same excuse, is almost as sensitive; it appears that he also was rather displeased with my article; he imagines that he alone knows the Kabbalah and is capable of speaking of it. He is in fear that Evola will soon do as much for the Tantras, although he is not very qualified to take it on; he sees all that through his philosophy, from where a deformation arises in the German manner. The true conception of Shakti is a completely different thing from `voluntarism'.
\end{quotex}


The second letter is from 26 January 1926

\begin{quotex}
Evola doesn't lack any pretentions, as you see; but, for my part, I continue to think that he does not understand at all what we mean by `intellectuality', `knowledge', `contemplation', etc., and that he doesn't even know how to make the distinction between the `initiatic' point of view and the `profane' point of view. It appears that he has the intention to publish a review of my work on the Vedanta in the journal Realistic Idealism; we will see what that will be. In any case, in spite of everything that we have tried to explain to him, he persists in finding `rationalism' in the Vedanta, all while failing to recognize that he then takes this word `rationalism' in a rather different sense that is usually given to it.
\end{quotex}



\flrightit{Posted on 2012-05-31 by Cologero}

\begin{center}* * *\end{center}

\begin{footnotesize}
\begin{sffamily}
\texttt{Zm on 2012-05-31 at 22:54 said:}

Yes, JE did gain wisdom in maturity.

What, pray, is the reason (if any) for the relative absence of references to F. Schuon here? His ``Observations'' on RG could be made the basis of some profitable discussion.

\hfill

\texttt{Cologero on 2012-05-31 at 23:39 said:}

Good point, Zm. Revolt, published nearly a decade later, did earn respect from Guenon despite his reservations. The simple answer to your question is that I am not as familiar with Schuon's works. In any case, there is abundant material available on Schuon elsewhere.

\hfill

\texttt{Abdul Mateen on 2012-06-07 at 07:30 said:}

I like your post, your post is very important for me.. contact me to visit \url{http://zulfiqarahmadnaqshbandi.com}

\hfill

\texttt{Numen on 2017-05-03 at 14:08 said:}

Quite illuminating, thank you for the translations, as always.

I remember reading a joke once about Evola and Guenon's relationship. Basically went like, Guenon was the only person who bought any of the French editions of Evola's work. It is interesting that, despite Evola's early arrogance and ignorance, Guenon seemed to see something in Evola worth correcting.

\end{sffamily}
\end{footnotesize}

